\scanpage{41}
\diagram{\begin{tikzcd}[column sep=small,row sep=large]
\operatorname{Hom}^d(\tau_\alpha M_{S_k},\underset{\substack{\cap\\\tau_{\alpha+d}N_{S_k}}}{\tau_\beta N_{S_k}})\arrow[r,hook] & \operatorname{Hom}^d(\tau_\alpha M_{S_k},\tau_{\alpha+d}N_{S_k})\arrow[d,phantom,"\wr"]\\
& \operatorname{Hom}_k(M,N)
\end{tikzcd}}
\margin{(sous-module des hom de degré $d$ de $\tau_\alpha M_{S_k}$ ds $\tau_{\alpha+d}N_{S_k}$ qui veulent bien appliquer $\tau_\alpha M_{S_k}$ dans $\tau_\beta N_{S_k}\subset\tau_{\alpha+d}N_{S_k}$)}

Il faut donc prendre un $u:M\longrightarrow N$ tel que l'hom correspondant $\tau_\alpha M_{S_k}-\tau_{\alpha+d}N_{S_k}$
\[
mT^{(j)}\longmapsto c_{dj}u(m)T^{(d+j)}
\]
\deleted{\unsure{applique} $M_S$} a son image dans $\tau_\beta\unsure{M|N}_{S_k}$, i.e. tel que
\[
\left\{\begin{array}{l}
c_{d,\alpha}u=0\\
c_{d,\alpha+1}u=0\\
\cdots\\
c_{d,\beta-d}u=0
\end{array}\right.
\]
(ce qui s'écrit
\[
c_{d,j}u=0\quad\text{si}\quad\alpha\leq j<\beta-d\quad\deleted{\unsure{\beta-d+1}}
\]
)

c) \uline{$d<0$}. Notons que
\[
\unsure{N^*_{\otimes S_k}|N^*_{S_k}}=\varphi_{0\beta}\tau_\beta N_{S_k}=\underset{\substack{\cap\\\varphi_\beta N_{S_k}}}{\tau_0\varphi_\beta\tau_\beta N_{S_k}}
\]
donc \inserted{si $M^*\in\operatorname{Ob}\mathcal M^*$}
\[
\begin{array}{ll}
\operatorname{Hom}^d(M^*,N_{S_k})\subset & \operatorname{Hom}^d(M^*,\varphi_\beta\tau_\beta N_{S_k})\\
& \Vert\\
& \operatorname{Hom}(M^{\unsure{*|\prime}}[-d],\varphi_\beta N_{S_k})\\
& \Vert\\
& \operatorname{Hom}(\tau_\beta M^{\unsure{*|\prime}}[-d],N_{S_k}
\end{array}
\]
\mspage{40}
